% Standalone solution dossier for the Route C exact cases.
\documentclass[../main.tex]{subfiles}
\begin{document}

% === SOLUTION HEADER (proof provenance) =====================================
% ledger-nodes: thm:cmh-1d; thm:cmh-product; cor:cmh-linear-images;
%               lem:cmh-gamma-completion; lem:cmh-row-min;
%               lem:cmh-angular-coefficient; thm:cmh-dirichlet;
%               cor:cmh-dirichlet-surplus; cor:cmh-dirichlet-poincare;
%               rem:cmh-saturation-risk
% refines: prog:cmh-route   (modules/kls/40-moment-map-cmh.tex)
% manuscript: modules/kls/42-cmh-exact-cases.tex
% depends: thm:cmh-implies-affine-poincare (solutions/thm-cmh-normalization.tex)
% bounded_by: none
%   The Eldan obstruction set of research/kls/obstructions.md is scoped in that
%   file to the fixed-cut program. obs:proj-ceiling is method-specific to radial /
%   projection tests and is not used here; see "Obstructions respected" below.
% evidence_run: none
%   Analytic. The finum target `cmh-gate-zero` regression-tests
%   eq:cmh-dirichlet-surplus on exact Dirichlet moments; it can refute but not prove.
% checked_by: agent
% authored_by: /root/kls_ledger_audit (repairs to the original dossier)
% reviewed_by: /root/cmh_exact_reviewer
% review: research/reviews/2026-08-25-kls-cmh-exact-cases-repair-audit.md
% date: 2026-08-25
% ============================================================================

\section*{Solution: exact CMH constants on the line, on products,
and on the log-concave Dirichlet family}
\label{sec:sol-cmh-dirichlet}

\paragraph{Scope.}
This dossier proves $\mathrm{CMH}(4)$ in the three classes where the constant of
Definition~\ref{def:cmh} is computable, the third being the first genuinely nonproduct family
for which Route~C has a theorem. Combined with
Theorem~\ref{thm:cmh-implies-affine-poincare} (dossier
\texttt{solutions/thm-cmh-normalization.tex}) each result yields an affine Poincar\'e bound with
the same constant. Nothing here bears on universal $\mathrm{CMH}(4)$, which remains open;
Corollary~\ref{rem:cmh-saturation-risk} in the manuscript explains why the product case in
particular is a warning rather than encouragement.

Notation is that of \S\ref{subsec:cmh-definition}: $H$ is the moment-map Stein kernel in target
coordinates, $L_\mu=\Div_\mu(H\nabla\,\cdot\,)$, $\Aop=-L_\mu$, and $\CMH$ is
\eqref{eq:cmh-constant}.

\subsection*{1. The line}

Let $\mu(\dd x)=\rho\dd x$ be centered on $(\ell,r)$ with variance $\sigma^2$ and Stein kernel
$\tau$, the zero-flux solution of $(\tau\rho)'=-x\rho$.

\begin{theorem}[= Theorem~\ref{thm:cmh-1d}]
\label{thm:sol-cmh-1d}
$\CMH(\mu)=\CP(\mu)/\Var(\mu)$. Every one-dimensional log-concave law satisfies
$\mathrm{CMH}(4)$, and the centered one-sided exponential has $\CMH=4$.
\end{theorem}

\begin{proof}
$L_\mu f=\tau f''-xf'=\rho^{-1}(\tau\rho f')'$, so for $h=-L_\mu f$ the field
$u=\tau f'$ satisfies $D_\mu^*u=h$ with vanishing flux at $\ell$, where $D=\dd/\dd x$ and
$D_\mu^*u=-\rho^{-1}(\rho u)'$. If $v=(D_\mu^*D)^{-1}h$, then $Dv$ solves the same equation with
the same boundary flux and hence equals $u$. Thus
\[
  \norm u_2^2=\norm{Dv}_2^2
  =\inner{v}{D_\mu^*Dv}
  =\inner{h}{(D_\mu^*D)^{-1}h},
\]
so $\sup_{h\perp1}\norm u_2^2/\norm h_2^2=\norm{(D_\mu^*D)^{-1}}=\CP(\mu)$, the last equality
because $D_\mu^*D$ is the nonnegative Langevin operator of $\mu$. In one dimension
$H=\tau$ and $\Sigma=\sigma^2$, so the CMH numerator is $\norm u_2^2/\sigma^2$ and the
denominator is $\norm h_2^2$; taking the supremum gives $\CMH=\CP/\sigma^2$.

The one-dimensional specialization of the Kannan--Lov\'asz--Simonovits bound
\cite{KannanLovaszSimonovits1995}, recorded for example in
\cite[Eq.~(2.25)]{cattiaux2018poincare} with that attribution, is
$\CP(\mu)\le4\Var(\mu)$.

For sharpness let $X=Y-1$ with $Y\sim\mathrm{Exp}(1)$. Then $X$ is centered, has density
$e^{-(x+1)}\one_{x\ge-1}\dd x$, and $\Var(X)=1$. For $0<a<\tfrac12$ set
$f_a(x)=e^{a(x+1)}$. Direct integration gives
\[
  \frac{\Var(f_a(X))}{\E\abs{f_a'(X)}^2}
  =\frac{(1-2a)^{-1}-(1-a)^{-2}}{a^2(1-2a)^{-1}}
  =\frac1{(1-a)^2}\longrightarrow4.
\]
The upper bound therefore gives $\CP(X)=4$, and the identity already proved gives $\CMH(X)=4$.
\end{proof}

\begin{remark}
This is an identity, not an inequality: in one dimension the solenoidal channel of
Proposition~\ref{prop:cmh-hodge} is empty (a square-integrable divergence-free field with
vanishing flux is zero), so CMH carries exactly the affine Poincar\'e information.
\end{remark}

\subsection*{2. Products}

\begin{theorem}[= Theorem~\ref{thm:cmh-product}]
\label{thm:sol-cmh-product}
For $\mu=\bigotimes_{i=1}^m\mu_i$ with centered factors for which the canonical CMH data are
defined,
\[
  \CMH(\mu)=\max_i\CMH(\mu_i).
\]
For one-dimensional factors this equals $\max_i\CP(\mu_i)/\Var(\mu_i)$. Consequently block
products of $\mathrm{CMH}(4)$ factors, and their invertible affine images, satisfy
$\mathrm{CMH}(4)$.
\end{theorem}

\begin{proof}
Here $\Sigma=\diag(\Sigma_1,\dots,\Sigma_m)$, $H=\diag(H_1,\dots,H_m)$, and
$L_\mu=\sum_iL_i$ with $L_i$ acting in block $i$ only, so the $L_i$ are commuting nonpositive
generators. For $i\ne j$, two integrations by parts give
\[
  \E[(L_ig)(L_jg)]
  =\E\Tr\!\left(H_iD^2_{ij}g\,H_jD^2_{ji}g\right)
  =\E\norm{H_i^{1/2}D^2_{ij}g\,H_j^{1/2}}_{\HS}^2\ge0.
\]
Hence
$\E(L_\mu g)^2=\sum_i\E(L_ig)^2+2\sum_{i<j}\E[(L_ig)(L_jg)]\ge\sum_i\E(L_ig)^2$.

For the numerator, conditioning on the other blocks and using the defining inequality for
$\CMH(\mu_i)$ gives
\[
  \E\inner{H_i\nabla_i g}{\Sigma_i^{-1}H_i\nabla_i g}
  \le\CMH(\mu_i)\,\E(L_i g)^2.
\]
Summing and using the previous display gives ``$\le$''. Testing on $g$ depending on one block
gives ``$\ge$''. The one-dimensional formula follows from Theorem~\ref{thm:sol-cmh-1d}; the
affine-image consequence uses the invertible affine covariance of $\CMH$ from
\S\ref{subsec:cmh-conventions}.
\end{proof}

\begin{corollary}[= Corollary~\ref{cor:cmh-linear-images}]
\label{cor:sol-cmh-linear-images}
If independent factors each satisfy the affine Poincar\'e inequality with constant $C$, so does
every linear image of their product, in particular every sum of independent such factors.
\end{corollary}

\begin{proof}
For $Y=TX$ and smooth $f$, $\Var_Yf=\Var_X(f\circ T)$ and
$\E_X\inner{\Cov(X)\nabla(f\circ T)}{\nabla(f\circ T)}
=\E\inner{T\Cov(X)T^\top\nabla f}{\nabla f}=\E\inner{\Cov(Y)\nabla f}{\nabla f}$. No
invertibility is needed.
\end{proof}

\subsection*{3. The Dirichlet family: setup}

Let $P\sim\Dir(\alpha)$, $\alpha_i\ge1$, $A=\sum_i\alpha_i$, $q=\alpha/A$. With
$C(p)=\diag(p)-pp^\top$ and $L_\alpha$ the Wright--Fisher generator
\eqref{eq:cmh-wright-fisher}, the potential
$\varphi(y)=A\log\sum_ie^{y_i/A}-q\cdot y$ has $\nabla\varphi=p-q$ and $D^2\varphi=C(p)/A$, and
pushes $e^{-\varphi}$ forward to the centered law $P-q$. Hence the canonical data are
$H=C(p)/A$, $L_\mu=L_\alpha/A$, and $\Sigma$ as in \eqref{eq:cmh-dirichlet-data}.

\begin{lemma}[Tangent pseudo-inverse]
\label{lem:sol-dirichlet-pseudoinverse}
For $v$ with $\sum_iv_i=0$, $v^\top\Sigma^\dagger v=A(A+1)\sum_iv_i^2/\alpha_i$.
\end{lemma}

\begin{proof}
Let $x=A(A+1)\diag(\alpha)^{-1}v$. Then
$\Sigma x=\bigl(\diag(\alpha)-\alpha\alpha^\top/A\bigr)\diag(\alpha)^{-1}v
=v-(\alpha/A)\sum_iv_i=v$. Since $\ker\Sigma=\mathrm{span}\{\one\}$ and $v^\top\one=0$, the
solution's ambiguity pairs to zero, so $v^\top\Sigma^\dagger v=v^\top x$.
\end{proof}

With $u_i=(C(P)\nabla g)_i$, $d_\alpha$, $n_\alpha$ as in \eqref{eq:cmh-dirichlet-dn}, the
statement $\CMH\le4$ becomes exactly $A(A+1)d_\alpha(g)\le4n_\alpha(g)$: the numerator is
$A^{-2}\cdot A(A+1)d_\alpha=\tfrac{A+1}Ad_\alpha$ and the denominator $A^{-2}n_\alpha$.

\begin{theorem}[= Theorem~\ref{thm:cmh-dirichlet}]
\label{thm:sol-cmh-dirichlet}
For every $m\ge2$, all $\alpha_i\ge1$, and all $g$ in the core,
$A(A+1)\,d_\alpha(g)\le4\,n_\alpha(g)$.
\end{theorem}

\subsection*{4. The Gamma lift}

Let $Y_i\sim\GammaLaw(\alpha_i,1)$ independent, $S=\sum_iY_i$, $P=Y/S$; then
$S\sim\GammaLaw(A,1)$, $P\sim\Dir(\alpha)$, and $S\perp P$. Set $G(Y)=g(Y/S)$, homogeneous of
degree $0$, and let $\calL_\Gamma G=\sum_i(Y_iG_{ii}+(\alpha_i-Y_i)G_i)$ be the Laguerre
generator, with
$N_\Gamma(G)=\E(\calL_\Gamma G)^2=\E[\sum_iY_iG_i^2+\sum_{i,j}Y_iY_jG_{ij}^2]$
its integrated Bochner identity, and $D_i(G)=\E[Y_i^2G_i^2]/\alpha_i$,
$D_\Gamma=\sum_iD_i$.

\begin{lemma}[= Lemma~\ref{lem:cmh-gamma-completion}]
\label{lem:sol-gamma-completion}
For $\alpha_i\ge1$, $N_\Gamma(G)-\tfrac14D_\Gamma(G)=\sum_i\E R_i+\sum_i\delta_iD_i(G)$ with
$R_i,\delta_i$ as in \eqref{eq:cmh-Ri-deltai}, both nonnegative.
\end{lemma}

\begin{proof}
For $Y\sim\GammaLaw(a,1)$ and smooth $u$ with the usual decay,
\[
  \E[Y^2uu']=\tfrac12\E[Y^2(u^2)']
  =\tfrac1{2\Gamma(a)}\int y^{a+1}e^{-y}(u^2)'\dd y
  =-\tfrac1{2\Gamma(a)}\int u^2\bigl((a+1)y^a-y^{a+1}\bigr)e^{-y}\dd y,
\]
i.e. $\E[Y^2uu']=-\tfrac12\E[((a+1)Y-Y^2)u^2]$. Expanding the first square of $R_i$ and applying
this with $u=G_i$, $u'=G_{ii}$, $a=\alpha_i$, the cross term is
\[
  -\tfrac2{\alpha_i+1}\E[Y_i^2G_iG_{ii}]
  =\tfrac1{\alpha_i+1}\E[((\alpha_i+1)Y_i-Y_i^2)G_i^2]
  =\E[Y_iG_i^2]-\tfrac1{\alpha_i+1}\E[Y_i^2G_i^2].
\]
Therefore
\[
  \E R_i=\E\Bigl[Y_iG_i^2+\sum_jY_iY_jG_{ij}^2\Bigr]
  +\Bigl(\tfrac1{(\alpha_i+1)^2}-\tfrac1{\alpha_i+1}\Bigr)\E[Y_i^2G_i^2],
\]
and $\tfrac1{(\alpha_i+1)^2}-\tfrac1{\alpha_i+1}=-\tfrac{\alpha_i}{(\alpha_i+1)^2}$. Summing over
$i$, the bracket sums to $N_\Gamma(G)$ and the correction to
$-\sum_i\tfrac{\alpha_i^2}{(\alpha_i+1)^2}D_i(G)$. Subtracting $\tfrac14D_\Gamma$ gives the
claim, and $\delta_i=\tfrac{\alpha_i^2}{(\alpha_i+1)^2}-\tfrac14\ge0$ iff
$\tfrac{\alpha_i}{\alpha_i+1}\ge\tfrac12$ iff $\alpha_i\ge1$.
\end{proof}

This lemma alone proves $\mathrm{CMH}(4)$ for the product-Gamma law; for that consequence,
$\alpha_i\ge1$ is used only to make $\delta_i\ge0$. In the Dirichlet proof below the exact
$\delta_i$ term is instead retained inside $F_A(\alpha_i,P_i)$, and log-concavity is consumed
when Lemma~\ref{lem:sol-angular} is applied with $a=\alpha_i\ge1$.

\begin{lemma}[= Lemma~\ref{lem:cmh-row-min}]
\label{lem:sol-row-min}
$R_i\ge\dfrac{Y_i}S\Bigl(1+\dfrac{Y_i}{\alpha_i+1}\Bigr)^2\abs{G_i}^2$.
\end{lemma}

\begin{proof}
Homogeneity of degree $0$ gives Euler's identity $\sum_jY_jG_j=0$; differentiating in $Y_i$
yields $G_i+\sum_jY_jG_{ji}=0$, i.e. the single linear constraint $\sum_jY_jG_{ij}=-G_i$. Set
$t_i=Y_i(G_{ii}-G_i/(\alpha_i+1))$ and $t_j=Y_jG_{ij}$ for $j\ne i$. Then
\[
  \sum_jt_j=\sum_jY_jG_{ij}-\tfrac{Y_i}{\alpha_i+1}G_i=-G_i\Bigl(1+\tfrac{Y_i}{\alpha_i+1}\Bigr),
  \qquad
  R_i=t_i^2+\sum_{j\ne i}\tfrac{Y_i}{Y_j}t_j^2 .
\]
The weights are $w_i=1$ and $w_j=Y_i/Y_j$, so $\sum_jw_j^{-1}=1+\sum_{j\ne i}Y_j/Y_i=S/Y_i$.
Minimizing $\sum_jw_jt_j^2$ subject to $\sum_jt_j=c$ gives $c^2/\sum_jw_j^{-1}$, which is the
stated bound.
\end{proof}

Homogeneity further gives $\calL_\Gamma G=S^{-1}L_\alpha g$ and $Y_iG_i=u_i(P)$: indeed
$G_i=\partial_{Y_i}g(Y/S)=S^{-1}(g_i-\sum_kP_kg_k)$, so $Y_iG_i=P_ig_i-P_i\sum_kP_kg_k=u_i(P)$.
For $A>2$, $\E S^{-1}=(A-1)^{-1}$ and $\E S^{-2}=z_A^{-1}$ with $z_A=(A-1)(A-2)$. Since
$S\perp P$,
\[
  N_\Gamma(G)=\E[S^{-2}]\,\E(L_\alpha g)^2=\frac{n_\alpha(g)}{z_A},
  \qquad
  D_\Gamma(G)=\E\sum_i\frac{(Y_iG_i)^2}{\alpha_i}=d_\alpha(g),
\]
the second because $Y_iG_i$ depends on $P$ alone. Substituting $Y_i=SP_i$ and $G_i=u_i/(SP_i)$
into Lemma~\ref{lem:sol-row-min} and taking expectations gives \eqref{eq:cmh-integrated-row}.

\subsection*{5. The scalar minimization}

\begin{lemma}[= Lemma~\ref{lem:cmh-angular-coefficient}]
\label{lem:sol-angular}
For $A\ge3$, $a\ge1$, $0<p\le1$ and $z=z_A$, the function $F_A$ of \eqref{eq:cmh-FA} satisfies
$F_A(a,p)\ge A-\tfrac12+s_A$ with $s_A$ as in \eqref{eq:cmh-FA-min}, and $s_A\ge0$.
\end{lemma}

\begin{proof}
The $p$-dependent part of $F_A$ is $\phi_a(p)=a\bigl(p^{-1}+zp(a+1)^{-2}\bigr)$, convex on
$(0,\infty)$ with unconstrained minimizer $p_*=(a+1)/\sqrt z$. Hence on $(0,1]$ the minimum is
at $p=1$ when $p_*\ge1$, i.e. $\sqrt z\le a+1$, and at $p_*$ otherwise.

\emph{Case $z\le4$.} Then $\sqrt z\le2\le a+1$ for every $a\ge1$, so $p=1$ and
\[
  F_A(a,1)=a+\frac{2a(A-2)}{a+1}+\frac{za}{a+1}-\frac z4 .
\]
Each of the three terms is increasing in $a$ on $[1,\infty)$, so the minimum is at $a=1$:
$1+(A-2)+z/2-z/4=A-1+z/4=A-\tfrac12+\tfrac{z-2}4$.

\emph{Case $z\ge4$.} At $p=p_*$ the two terms of $\phi_a$ are equal, so
$\phi_a(p_*)=2a\sqrt z/(a+1)$. Writing $r=a/(a+1)\in[\tfrac12,1)$ for $a\ge1$,
\[
  F_A=2r\sqrt z+2r(A-2)+zr^2-\frac z4=2r(\sqrt z+A-2)+zr^2-\frac z4,
\]
increasing in $r$ on $[0,1)$, so the minimum is at $r=\tfrac12$ ($a=1$) and equals
$(\sqrt z+A-2)+z/4-z/4=A-2+\sqrt z=A-\tfrac12+\sqrt z-\tfrac32$. In the complementary region
$\sqrt z\le a+1$ the value $F_A(a,1)$ is increasing in $a$ and at the interface is not smaller,
so the stated minimum stands.

Finally $A\ge3$ gives $z_A=(A-1)(A-2)\ge2$, so $s_A=(z-2)/4\ge0$ when $z\le4$ and
$s_A=\sqrt z-\tfrac32\ge2-\tfrac32>0$ when $z\ge4$.
\end{proof}

\subsection*{6. Proof of the Dirichlet theorem}

\begin{proof}[Proof of Theorem~\ref{thm:sol-cmh-dirichlet}]
If $m=2$ and $A<3$, the law is one dimensional and Theorem~\ref{thm:sol-cmh-1d} applies (a
$\mathrm{Beta}$ law is log-concave precisely when both parameters are $\ge1$). In every
remaining case $A=\sum_i\alpha_i\ge3$ --- automatically so when $m\ge3$ --- hence $z_A\ge2>0$
and the inverse moments $\E S^{-1},\E S^{-2}$ are finite.

By Lemma~\ref{lem:sol-gamma-completion} and \eqref{eq:cmh-integrated-row}, since
$D_i(G)=\E[u_i^2]/\alpha_i$, the coefficient of $\E_P[u_i^2]$ in
$N_\Gamma(G)-\tfrac14D_\Gamma(G)$ is at least
\[
  \frac1{z_AP_i}+\frac2{(A-1)(\alpha_i+1)}+\frac{P_i}{(\alpha_i+1)^2}+\frac{\delta_i}{\alpha_i}.
\]
Multiply by $z_A\alpha_i$ and use $z_A/(A-1)=A-2$ together with
$z_A\delta_i=z_A\alpha_i^2/(\alpha_i+1)^2-z_A/4$:
\[
  \frac{\alpha_i}{P_i}+\frac{2\alpha_i(A-2)}{\alpha_i+1}
  +\frac{z_A\alpha_i(\alpha_i+P_i)}{(\alpha_i+1)^2}-\frac{z_A}4
  =F_A(\alpha_i,P_i).
\]
Lemma~\ref{lem:sol-angular} bounds this below by $A-\tfrac12+s_A$ pointwise in $P_i$, so
\[
  N_\Gamma(G)-\tfrac14D_\Gamma(G)
  \ \ge\ \frac{A-\tfrac12+s_A}{z_A}\sum_i\frac{\E_P[u_i^2]}{\alpha_i}
  =\frac{A-\tfrac12+s_A}{z_A}\,d_\alpha(g).
\]
Substituting $N_\Gamma=n_\alpha/z_A$ and $D_\Gamma=d_\alpha$, multiplying by $z_A$, and using
\[
  \frac{z_A}4+A-\frac12=\frac{A^2-3A+2}4+\frac{4A-2}4=\frac{A(A+1)}4,
\]
we obtain $n_\alpha(g)\ge\bigl(\tfrac{A(A+1)}4+s_A\bigr)d_\alpha(g)$. Since $s_A\ge0$ this gives
$A(A+1)d_\alpha(g)\le4n_\alpha(g)$.
\end{proof}

\begin{corollary}[= Corollary~\ref{cor:cmh-dirichlet-surplus}]
For $A>3$, $\CMH\le4\bigl(1+4s_A/(A(A+1))\bigr)^{-1}<4$.
\end{corollary}

\begin{proof}
For $A>3$ the Gamma branch of the preceding proof applies, including when $m=2$, and gives
$n_\alpha\ge\bigl(A(A+1)/4+s_A\bigr)d_\alpha$. Since $z_A>2$, the definition of $s_A$ gives
$s_A>0$, and rearrangement proves the claim.
\end{proof}

\begin{corollary}[= Corollary~\ref{cor:cmh-dirichlet-poincare}]
Every log-concave Dirichlet law has $\CPaff\le4$, as does every product, linear image, and
convolution of independent such laws.
\end{corollary}

\begin{proof}
Theorem~\ref{thm:cmh-implies-affine-poincare} applied to
Theorem~\ref{thm:sol-cmh-dirichlet}, then Corollary~\ref{cor:sol-cmh-linear-images}.
\end{proof}

Qualitative dimension-free KLS bounds for simplices and conservative Gamma models are prior
art; see \cite[\S1.2]{KolesnikovMilman2016OrliczKLS} as a secondary pointer. The specific
content of this proof is the constant $4$, the quantitative surplus, and the CMH-level estimate;
no broader priority claim is made.

\paragraph{Closure.}
The argument is written for smooth $g$ with controlled boundary behavior. To close: take
polynomials on the simplex, lift them after a radial cutoff $\{S\ge\eps\}$ so that all Gamma
integrations by parts in Lemma~\ref{lem:sol-gamma-completion} are justified, and let
$\eps\downarrow0$; the inverse moments $\E S^{-1},\E S^{-2}$ used in
\eqref{eq:cmh-radial-relations} are finite because $A\ge3$ throughout the Gamma branch.
Polynomials form a core for $L_\alpha$, so the closed forms extend the inequality to
$\Dom(L_\alpha)$. Only the residual branch $m=2$, $A<3$ uses the one-dimensional no-flux closure
of Theorem~\ref{thm:sol-cmh-1d} instead.

\paragraph{Obstructions respected.}
No \texttt{bounded\_by} edge applies: the entries of \texttt{research/kls/obstructions.md} are
scoped there to the fixed-cut Eldan program. The only one with method-level reach,
\texttt{obs:proj-ceiling}, forbids deriving quadratic-chaos thin shell from radial or projection
information alone; the Dirichlet proof uses the full Hessian row through the Euler constraint of
Lemma~\ref{lem:sol-row-min}, not projection tests, and makes no thin-shell claim.
Consistency with the sharp external inputs holds: by Theorem~\ref{thm:cmh-1d}, the boundary
mechanism forcing the constant $4$ is the one-dimensional exponential, which is also the
extremal case of Theorem~\ref{thm:letwin-moment-map}.

\begin{corollary}[= Corollary~\ref{rem:cmh-saturation-risk}]
\label{rem:sol-cmh-saturation-risk}
Theorems~\ref{thm:sol-cmh-1d} and \ref{thm:sol-cmh-product} prove that products of centered
one-sided exponentials have $\CMH=4$ exactly. Proposition~\ref{prop:cmh-hodge} proves, for each
admissible test function, the exact decomposition of the CMH numerator into the affine
Poincar\'e contribution and a nonnegative solenoidal contribution.
\end{corollary}

\begin{remark}[Perturbative caveat]
It is an open question, not a consequence of this corollary, whether an admissible perturbation
raises the full CMH Rayleigh quotient. The covariance, canonical Stein kernel, both numerator
channels, denominator, and optimizer all vary, so increasing a solenoidal term in isolation is
insufficient. The required full second-variation calculation is
Question~\ref{q:cmh-solenoidal-perturbation}.
\end{remark}

\paragraph{What this does not show.}
Theorem~\ref{thm:sol-cmh-dirichlet} is a family result, not evidence for universal
$\mathrm{CMH}(4)$. By Corollary~\ref{cor:cmh-dirichlet-surplus} the Dirichlet family is
strictly inside the bound, and by Theorem~\ref{thm:sol-cmh-product} the saturating cases are
products of centered one-sided exponentials, where the slack is exactly zero. Whether that
endpoint is perturbatively unstable is precisely the open issue just stated.
\end{document}
